% !TeX root = ../tesis.tex

\section{Experimentando con hashear certificados}

Durante una etapa intermedia del prototipo se exploró una estrategia conceptualmente sencilla para reforzar el \textit{binding} entre el estado publicado \Onchain y la evidencia criptográfica asociada: recomputar dentro del validador el hash completo de los certificados de \Mithril y verificar que dicho hash coincidiera con el valor almacenado en el propio certificado. En el código de esa etapa, esta idea aparece encapsulada en la función \texttt{verify\_hash\_matches\_content}, que compara el campo \texttt{certificate.hash} con el resultado de volver a ejecutar \texttt{hash\_certificate(certificate)} en Aiken.

La motivación inicial era razonable. Si el validador puede rehashear el certificado completo, entonces obtiene una forma muy directa de comprobar que el objeto recibido en el datum o en el redeemer coincide exactamente con la estructura que luego se usa para encadenar certificados, verificar \textit{epoch chaining} y validar mensajes firmados. Desde un punto de vista puramente lógico, la estrategia parecía atractiva porque reducía la cantidad de supuestos auxiliares: el propio script reconstruía el compromiso criptográfico del certificado sin depender de una abstracción adicional.

Sin embargo, el experimento mostró rápidamente que la idea era demasiado costosa para el modelo de ejecución de \Cardano. El problema no estaba solamente en aplicar \texttt{sha2\_256}, sino en toda la serialización y transformación intermedia que el proceso exige. En la implementación estudiada, el rehash de un certificado obliga a procesar varios subcomponentes, incluyendo metadatos, mensajes de protocolo, claves agregadas y firmas. Además, la representación heredada del cliente de \Mithril hacía que algunos campos se manipularan como cadenas hexadecimales, duplicando de hecho el volumen de bytes arrastrado por ciertas operaciones. El resultado fue una presión excesiva tanto sobre memoria como sobre CPU, incluso antes de sumar el resto de la lógica del validador.

Para cuantificar ese costo, se reutilizaron los tests unitarios del proyecto Aiken y se ejecutaron con \texttt{aiken check}, tomando los \textit{execution units} reportados por la herramienta para cada caso \cite{aiken2026}. Aunque estas mediciones no reemplazan un benchmark final \Onchain con transacciones completas, sí ofrecen una aproximación muy útil porque reflejan el costo UPLC de los programas efectivamente generados por el compilador. La Tabla~\ref{tab:chapter9-hash-benchmarks} resume los casos más representativos.

\begin{table}[ht]
  \centering
  \small
  \begin{tabular}{p{6.0cm}r r}
    \toprule
    \textbf{Test de Aiken} & \textbf{Mem} & \textbf{CPU} \\
    \midrule
    \texttt{hash\_certificate\_metadata\_test} & 1,857,590 & 531,863,069 \\
    \texttt{hash\_certificate\_test} & 41,725,138 & 11,696,733,751 \\
    \texttt{hash\_real\_cardano\_transactions}\\\texttt{\_certificate\_test} & 35,621,336 & 10,146,407,782 \\
    \texttt{hash\_real\_sd\_certificate\_test} & 63,613,144 & 17,634,279,602 \\
    \texttt{minting\_validator\_test} & 91,275,468 & 29,934,916,764 \\
    \texttt{stake\_distribution\_validator}\\\texttt{\_sd\_certificate\_test} & 129,762,816 & 37,939,818,729 \\
    \bottomrule
  \end{tabular}
  \caption{Mediciones de \textit{execution units} obtenidas con \texttt{aiken check} sobre pruebas del prototipo intermedio.}
  \label{tab:chapter9-hash-benchmarks}
\end{table}

Las cifras dejan ver dos hechos importantes. En primer lugar, hashear solamente los metadatos del certificado resulta relativamente barato en comparación con el resto del pipeline: el test \texttt{hash\_certificate\_metadata\_test} queda en el orden de millones bajos de memoria y cientos de millones de CPU. En segundo lugar, el salto al hash completo del certificado cambia drásticamente la escala del costo. El caso \texttt{hash\_real\_sd\_certificate\_test}, que rehashea un certificado real de \StakeDistribution, alcanza 63,613,144 unidades de memoria y 17,634,279,602 de CPU. Es decir, una porción muy significativa del presupuesto de una validación compleja se consumía antes siquiera de verificar la prueba criptográfica principal.

Este resultado se vuelve todavía más claro cuando se lo compara con tests de validadores completos. El test \texttt{minting\_validator\_test} consume 91,275,468 de memoria y 29,934,916,764 de CPU; por lo tanto, sólo el rehash de un certificado real de \StakeDistribution representa una fracción sustancial del costo total de una transacción de \textit{minting} del prototipo intermedio. En el caso de \texttt{stake\_distribution\_validator\_sd\_certificate\_test}, el presupuesto total sube a 129,762,816 de memoria y 37,939,818,729 de CPU, lo que confirma que la ruta basada en rehashear certificados completos dentro del script tensiona fuertemente los límites \Onchain aun antes de incorporar circuitos más elaborados.

La conclusión de ingeniería fue, entonces, no conservar esta estrategia en la implementación final. El objetivo de \textit{binding} seguía siendo correcto, pero no era necesario alcanzarlo recomputando el hash completo del certificado en Aiken. En su lugar, la versión final adopta un criterio más económico: publicar el certificado relevante como parte del estado autenticado, enlazarlo con entradas públicas bien elegidas en los circuitos y verificar \Onchain solamente las relaciones estrictamente necesarias para garantizar unicidad, encadenamiento y coherencia entre datum, redeemer y prueba. Dicho de otro modo, el sistema preserva el mismo objetivo semántico, pero desplaza la parte costosa del trabajo hacia circuitos especializados donde el hashing resulta mucho más natural.

Esta decisión también mejora la modularidad del sistema. Al sacar del validador el rehash completo del certificado, el script deja de depender de detalles accidentales de serialización y de conversión entre representaciones de bytes y cadenas. Esto reduce el acoplamiento con una forma particular de modelar los certificados de \Mithril en Aiken y deja una interfaz más limpia entre la capa de actualización y la capa \ZK. En retrospectiva, el experimento fue valioso precisamente porque permitió medir con rapidez una solución que, aunque correcta desde el punto de vista lógico, no era sostenible dentro de los presupuestos del protocolo.
